\section{Related Work}
\label{sec:related}

\subsection{O-RAN Simulation Platforms}
Lacava \emph{et~al.}~\cite{ns3_oran2021} introduced ns-O-RAN, coupling NS-3
with an O-DU/O-CU emulator and a near-RT RIC to study traffic steering via
E2 interface reports.
Our work extends this by adding 6G PHY abstractions (THz, massive MIMO) and
an integrated MEC service layer.

\subsection{AI/ML in O-RAN}
Reinforcement learning has been applied to handover decisions in
LTE/NR~\cite{dqn_handover2020}, showing that model-free agents can
outperform threshold-based rules under non-stationary traffic.
Federated learning~\cite{fedavg2017} allows gNBs to collaboratively train
a shared model without centralising raw measurements, preserving data
locality---a critical requirement in O-RAN deployments.

\subsection{Digital Twins for Networks}
Wu \emph{et~al.}~\cite{digital_twin_net2022} survey digital-twin
architectures for IoT and 5G, highlighting predictive accuracy and
synchronisation delay as key metrics.
We adopt a lightweight regression-based twin that updates at each simulation
step and is used to pre-validate RL actions.

\subsection{MEC Integration}
The ETSI MEC framework~\cite{mec_etsi2022} defines APIs for application
lifecycle management at the network edge.
Our simulation models MEC hosts co-located with gNBs and evaluates CPU/memory
utilisation and end-to-end service latency under varying UE densities.

\subsection{6G Candidate Technologies}
THz communications offer multi-Tbps capacities at the cost of severe
distance-dependent path loss~\cite{thz_6g2022}.
We model THz links between UEs and small cells and study their interaction
with MEC offloading decisions.

% TODO: expand with 3–5 more directly competing works once paper
%       submission venue is chosen.
